\section{Q1: Detection}
\label{sec:q1}

The formal matrices contain 448 bug images, each paired with its negative control: 352 for ASan and
Memcheck over all seven layers, and 96 for HWASan over \code{malloc}, Arena and Free list. All
positive controls were detected, all 448 negative controls were clean, and no image ended in a tool
error. Table~\ref{tab:summary} and Figure~\ref{fig:summary} summarise the verdicts;
Tables~\ref{tab:matrix-core} and~\ref{tab:matrix-sqlite} give every cell. A fraction counts
\textsc{detected} over eligible images, i.e.\ excluding \textsc{unsupported} ones.

\begin{table}[t]
\centering\small
\begin{tabular}{l*{3}{cc}}
\toprule
 & \multicolumn{2}{c}{ASan} & \multicolumn{2}{c}{Memcheck} & \multicolumn{2}{c}{HWASan} \\
\cmidrule(lr){2-3}\cmidrule(lr){4-5}\cmidrule(lr){6-7}
allocator layer & plain & annotated & plain & annotated & plain & annotated \\
\midrule
malloc baseline & 12/13 & n/a & 13/13 & n/a & 13/15 & n/a \\
Arena (ours) & 1/14 & 10/14 & 2/14 & 13/14 & 2/14 & 9/14 \\
Free list (ours) & 1/16 & 11/16 & 2/16 & 14/16 & 2/16 & 11/16 \\
SQLite memsys5 & 1/16 & 11/16 & 2/16 & 14/16 & -- & -- \\
SQLite lookaside & 1/16 & n/a & 2/16 & n/a & -- & -- \\
Arena rewrite & 12/13 & n/a & 13/13 & n/a & -- & -- \\
Free list rewrite & 12/13 & n/a & 13/13 & n/a & -- & -- \\
\bottomrule
\end{tabular}

\caption{Bugs detected out of eligible bugs, per allocator layer, without (plain) and with the tool's
annotation. ``n/a'': the layer has no annotated variant; ``--'': not run. The \code{malloc} layer was
also built with annotation on and gave identical results for ASan and Memcheck; HWASan detected case
4 on that build (\S\ref{sec:q1-precision}).}
\label{tab:summary}
\end{table}

\begin{figure}[t]
\centering
\begin{tikzpicture}
\begin{groupplot}[
  group style={group size=3 by 1, horizontal sep=0.9cm, yticklabels at=edge left},
  width=0.36\textwidth, height=4.6cm,
  ybar=1.5pt, /pgf/bar width=5.5pt,
  ymin=0, ymax=105, ytick={0,25,50,75,100},
  ymajorgrids, grid style={black!12}, axis line style={black!60}, tick style={black!60},
  symbolic x coords={malloc,arena,freelist,memsys5,lookaside,rewrite},
  xtick=data, xticklabel style={rotate=40, anchor=east, font=\scriptsize},
  yticklabel style={font=\scriptsize}, title style={font=\small},
  enlarge x limits=0.12,
  unbounded coords=discard,
  legend style={font=\scriptsize, draw=none, fill=none, legend columns=2, at={(0.5,1.30)}, anchor=south,
    /tikz/every even column/.append style={column sep=8pt}},
]
\nextgroupplot[title=ASan, ylabel={\small detected / eligible (\%)},
  xticklabels={malloc,Arena,Free list,memsys5,lookaside,rewrite}]
\addplot[fill=seriesA, draw=none] table[x=layer, y=plain] {tables/detect-asan.dat};
\addplot[fill=seriesB, draw=none] table[x=layer, y=annotated] {tables/detect-asan.dat};
\nextgroupplot[title=Valgrind Memcheck, xticklabels={malloc,Arena,Free list,memsys5,lookaside,rewrite}]
\addplot[fill=seriesA, draw=none] table[x=layer, y=plain] {tables/detect-valgrind.dat};
\addplot[fill=seriesB, draw=none] table[x=layer, y=annotated] {tables/detect-valgrind.dat};
\legend{unannotated, annotated}
\nextgroupplot[title=HWASan (AArch64), symbolic x coords={malloc,arena,freelist},
  xticklabels={malloc,Arena,Free list}, enlarge x limits=0.3]
\addplot[fill=seriesA, draw=none] table[x=layer, y=plain] {tables/detect-hwasan.dat};
\addplot[fill=seriesB, draw=none] table[x=layer, y=annotated] {tables/detect-hwasan.dat};
\end{groupplot}
\end{tikzpicture}
\caption{Share of eligible bugs detected per allocator layer (data of Table~\ref{tab:summary}).
\code{malloc}, lookaside and the rewrites have no annotated variant; ``rewrite'' is the Free list
rewrite (the Arena rewrite has the same verdicts).}
\label{fig:summary}
\end{figure}

\begin{table}[p]
\centering\footnotesize
\setlength{\tabcolsep}{3.4pt}
\renewcommand{\arraystretch}{1.08}
\begin{tabular}{r l|ccccc|ccccc|ccccc}
\toprule
 & & \multicolumn{5}{c}{ASan} & \multicolumn{5}{c}{Memcheck} & \multicolumn{5}{c}{HWASan} \\
 & case & \rot{malloc} & \rot{Arena} & \rot{Arena+A} & \rot{Free list} & \rot{Free list+A} & \rot{malloc} & \rot{Arena} & \rot{Arena+A} & \rot{Free list} & \rot{Free list+A} & \rot{malloc} & \rot{Arena} & \rot{Arena+A} & \rot{Free list} & \rot{Free list+A} \\
\midrule
0 & overflow, write +1 & \cD & \cM & \cD & \cM & \cD & \cD & \cM & \cD & \cM & \cD & \cD & \cM & \cM & \cM & \cM \\
1 & overflow, read +1 & \cD & \cM & \cD & \cM & \cD & \cD & \cM & \cD & \cM & \cD & \cD & \cM & \cM & \cM & \cM \\
2 & overflow into slack & \cD & \cM & \cD & \cM & \cD & \cD & \cM & \cD & \cM & \cD & \cD & \cM & \cM & \cM & \cM \\
3 & overflow into live obj. & \cM & \cM & \cM & \cM & \cM & \cD & \cM & \cM & \cM & \cM & \cD & \cM & \cD & \cM & \cD \\
4 & read past the pool & \cD & \cM & \cMs & \cMs & \cM & \cD & \cD & \cD & \cD & \cD & \cM & \cD & \cD & \cD & \cD \\
5 & underflow, write -1 & \cD & \cM & \cD & \cM & \cD & \cD & \cM & \cD & \cM & \cD & \cD & \cM & \cD & \cM & \cD \\
6 & underflow, read -1 & \cD & \cM & \cD & \cM & \cD & \cD & \cM & \cD & \cM & \cD & \cD & \cM & \cD & \cM & \cD \\
7 & underflow into header & \cD & \cM & \cD & \cM & \cD & \cD & \cM & \cD & \cM & \cD & \cD & \cM & \cD & \cM & \cD \\
8 & UAF write & \cD & \cM & \cD & \cM & \cD & \cD & \cM & \cD & \cM & \cD & \cD & \cM & \cD & \cM & \cD \\
9 & UAF read & \cD & \cM & \cD & \cM & \cD & \cD & \cM & \cD & \cM & \cD & \cD & \cM & \cD & \cM & \cD \\
10 & UAF after reuse & \cU & \cU & \cU & \cM & \cM & \cU & \cU & \cU & \cM & \cM & \cD & \cU & \cU & \cM & \cD \\
11 & UAF after realloc & \cD & \cM & \cD & \cM & \cD & \cD & \cM & \cD & \cM & \cD & \cD & \cM & \cD & \cM & \cD \\
12 & UAF after pool destroy & \cU & \cD & \cD & \cD & \cD & \cU & \cD & \cD & \cD & \cD & \cU & \cD & \cD & \cD & \cD \\
13 & double free & \cD & \cM & \cM & \cM & \cM & \cD & \cM & \cD & \cM & \cD & \cD & \cM & \cM & \cM & \cM \\
14 & interior free & \cD & \cM & \cM & \cM & \cM & \cD & \cM & \cD & \cM & \cD & \cM & \cM & \cM & \cM & \cM \\
15 & double free after reuse & \cU & \cU & \cU & \cM & \cD & \cU & \cU & \cU & \cM & \cD & \cD & \cU & \cU & \cM & \cD \\
\midrule
 & detected / eligible & \scriptsize 12/13 & \scriptsize 1/14 & \scriptsize 10/14 & \scriptsize 1/16 & \scriptsize 11/16 & \scriptsize 13/13 & \scriptsize 2/14 & \scriptsize 13/14 & \scriptsize 2/16 & \scriptsize 14/16 & \scriptsize 13/15 & \scriptsize 2/14 & \scriptsize 9/14 & \scriptsize 2/16 & \scriptsize 11/16 \\
\bottomrule
\end{tabular}

\caption{Detection matrix for the core layers (D = detected, M = missed, M$^{*}$ = missed, the
program crashed inside the tool's signal handler, U = unsupported, +A = annotated). ASan and Memcheck:
lab server; HWASan: AArch64 VM. The \code{malloc}+A column is omitted (identical to \code{malloc},
except HWASan case 4).}
\label{tab:matrix-core}
\vspace{1.2em}
\begin{tabular}{r l|ccccc|ccccc}
\toprule
 & & \multicolumn{5}{c}{ASan} & \multicolumn{5}{c}{Memcheck} \\
 & case & \rot{memsys5} & \rot{memsys5+A} & \rot{lookaside} & \rot{Arena rewr.} & \rot{Free-list rewr.} & \rot{memsys5} & \rot{memsys5+A} & \rot{lookaside} & \rot{Arena rewr.} & \rot{Free-list rewr.} \\
\midrule
0 & overflow, write +1 & \cM & \cD & \cM & \cD & \cD & \cM & \cD & \cM & \cD & \cD \\
1 & overflow, read +1 & \cM & \cD & \cM & \cD & \cD & \cM & \cD & \cM & \cD & \cD \\
2 & overflow into slack & \cM & \cD & \cM & \cD & \cD & \cM & \cD & \cM & \cD & \cD \\
3 & overflow into live obj. & \cM & \cM & \cM & \cM & \cM & \cM & \cM & \cM & \cD & \cD \\
4 & read past the pool & \cMs & \cMs & \cM & \cD & \cD & \cD & \cD & \cD & \cD & \cD \\
5 & underflow, write -1 & \cM & \cD & \cM & \cD & \cD & \cM & \cD & \cM & \cD & \cD \\
6 & underflow, read -1 & \cM & \cD & \cM & \cD & \cD & \cM & \cD & \cM & \cD & \cD \\
7 & underflow into header & \cM & \cD & \cM & \cD & \cD & \cM & \cD & \cM & \cD & \cD \\
8 & UAF write & \cM & \cD & \cM & \cD & \cD & \cM & \cD & \cM & \cD & \cD \\
9 & UAF read & \cM & \cD & \cM & \cD & \cD & \cM & \cD & \cM & \cD & \cD \\
10 & UAF after reuse & \cM & \cM & \cM & \cU & \cU & \cM & \cM & \cM & \cU & \cU \\
11 & UAF after realloc & \cM & \cD & \cM & \cD & \cD & \cM & \cD & \cM & \cD & \cD \\
12 & UAF after pool destroy & \cD & \cD & \cD & \cU & \cU & \cD & \cD & \cD & \cU & \cU \\
13 & double free & \cM & \cM & \cM & \cD & \cD & \cM & \cD & \cM & \cD & \cD \\
14 & interior free & \cM & \cM & \cM & \cD & \cD & \cM & \cD & \cM & \cD & \cD \\
15 & double free after reuse & \cM & \cD & \cM & \cU & \cU & \cM & \cD & \cM & \cU & \cU \\
\midrule
 & detected / eligible & \scriptsize 1/16 & \scriptsize 11/16 & \scriptsize 1/16 & \scriptsize 12/13 & \scriptsize 12/13 & \scriptsize 2/16 & \scriptsize 14/16 & \scriptsize 2/16 & \scriptsize 13/13 & \scriptsize 13/13 \\
\bottomrule
\end{tabular}

\caption{Detection matrix for SQLite's allocators and the two rewrites (lab server).}
\label{tab:matrix-sqlite}
\end{table}

\subsection{Unannotated CMAs hide most object-level errors}
\label{sec:q1-blind}

On \code{malloc} ASan detected 12 of 13 eligible bugs, Memcheck 13 of 13 and HWASan 13 of 15. In
every unannotated custom layer each tool detected one or two. The cells it kept are telling: case 12
(the pool itself is freed, so the one block the tool knows ends) for all three tools, and case 4 (the
read leaves the pool) for Memcheck and HWASan. Both are detections of the \emph{pool}, not of the
object, and should not be credited as allocator awareness. Everything that stays inside the pool ---
overflows into slack or into a neighbour, underflows, uses after free, invalid and double frees --- is
missed by every tool on Arena, Free list, memsys5 and lookaside alike. SQLite's allocators behave
exactly like our own: one of 16 for ASan and two of 16 for Memcheck. These 122 misses are the
largest class in the review of \S\ref{sec:audit}.

\subsection{Annotation recovers most bugs but leaves mechanism-specific gaps}
\label{sec:q1-annot}

With annotation, ASan detected 10 of 14 bugs on Arena and 11 of 16 on Free list and memsys5;
Memcheck 13 of 14 and 14 of 16; HWASan 9 of 14 and 11 of 16. The bugs that remain show that each
interface describes a different thing about memory, and that is what bounds it.

\paragraph{Poisoning describes addressability (ASan).} An annotated allocator can poison slack and
freed objects, so overflows by one byte, underflows and uses after free are caught. Three things
cannot be expressed. An overflow that lands in another \emph{live} object (case 3) touches bytes that
must stay unpoisoned; ASan missed case 3 on every layer, including \code{malloc}, where offset 64 clears
its 16-byte redzone and lands in the next block. Once the Free list reissues a slot (case 10), the
stale pointer and the new owner's pointer hold one address, and that address is legitimately
addressable. An invalid or double \code{free} (cases 13, 14) need not cause an invalid memory access,
even if the allocator reads or writes metadata. Poisoning alone does not track ownership, so ASan
missed both on every annotated layer although its own \code{malloc} reports them.

\paragraph{Memory pools describe ownership (Memcheck).} \code{MEMPOOL\_FREE} is checked against the
set of live pool objects, so Memcheck reports double and interior frees inside a pool, which
poisoning alone does not identify. It is still address-based, so it shares case 10 and case 3 with ASan: on
annotated Free list and memsys5 it missed exactly those two. On \code{malloc} it does detect case 3,
because its own allocator places a redzone on each side of every block and offset 64 falls between
blocks; our original prediction (missed) was wrong.

\paragraph{Tags describe provenance (HWASan).} A tag records which allocation a pointer belongs to,
so annotated HWASan detects what poisoning cannot: case 3, where the neighbour carries a different
tag, and case 10, where the reissued slot was given a new tag that the stale pointer lacks. Its own
limits are different. The tagging interface works on whole 16-byte granules and, unlike HWASan's
\code{malloc}, cannot describe a short last granule, so a 24-byte object's bytes 24--31 keep its tag
and the one-byte overflows of cases 0--2 are missed. Like poisoning, tags have no event for a bad
\code{free} (cases 13, 14).

\begin{table}[t]
\centering\small
\begin{tabular}{lccc}
\toprule
what the bug needs the tool to know & ASan poisoning & Memcheck pools & HWASan tags \\
\midrule
requested size, to the byte (cases 0--2) & yes & yes & no (16-byte granules) \\
object lifetime, before reuse (8, 9, 11) & yes & yes & yes \\
which object an address belongs to (3, 10) & no & no & yes \\
ownership at \code{free} (13, 14) & no & yes & no \\
\bottomrule
\end{tabular}
\caption{What each annotation interface can express, as measured on the annotated Arena and Free list.}
\label{tab:expressive}
\end{table}

Table~\ref{tab:expressive} condenses these results. None of the three interfaces covers every case,
and no two have the same blind spots, so the choice of tool determines which bugs inside an annotated
CMA can be found at all.

\subsection{Leaving the pool, and precision effects}
\label{sec:q1-precision}

Case 4 reads 8192 bytes past the end of the 1~MiB pool. We expected every tool to see it, since the
access leaves the only block the tool knows. Memcheck detects it on all tested layers, and HWASan
detects it on the tested pools. ASan does not, on
any custom layer, with or without annotation: a request of 1~MiB is served by ASan's large-allocation
path as a separate mapping whose redzone ends with its last page, so 8192 bytes further on lies
outside anything ASan has poisoned. The read is then either silent, if another mapping happens to be
there, or a segmentation fault caught by ASan's signal handler, if nothing is. Which of the two
happens depends on the address-space layout: on EC2 all four Arena and Free list cells crashed, on the
lab server two did. We record both as missed. This is the one place where a CMA makes ASan
\emph{worse} than Memcheck at a bug outside every object, and it is the same shape as the real SQLite
over-reads in \S\ref{sec:tier2}.

HWASan's case 4 on \code{malloc} was missed in the formal run although the case reduces to a one-byte
overflow there. Re-running that image 200 times gave 198 detections. The miss is consistent with
the random pointer tag equalling the short-granule value, allowing the check to pass; the reruns
support this explanation but do not establish the original run's tag values. We kept the measured
verdict rather than replacing it with a rerun.

\subsection{Rewriting the CMA away}
\label{sec:q1-rewrite}

The Arena and Free list rewrites detected 12 of 13 bugs under ASan and 13 of 13 under Memcheck ---
exactly \code{malloc}'s numbers --- with 14 lines of code each. Rewriting changes allocator semantics. Under
it, freed storage is not reissued immediately and there is no pool, so cases 10, 12 and 15 become
\textsc{unsupported}: the bug they model no longer exists. A rewrite also drops whatever the CMA was
for --- bulk deallocation, a memory bound, speed (\S\ref{sec:perf}) --- and changes allocation
order, sizes and reuse. A detection under the rewrite shows that the tool sees a libc object, not that
it understands the original allocator. For SQLite the rewrite needs no code at all: compiling without
memsys5 and with lookaside and the page-cache bulk pool disabled is the transformation we use as
the positive control of the real-bug check below.

\subsection{Reproducibility and review of the misses}
\label{sec:audit}

The ASan and Memcheck matrices were first produced on EC2 from an earlier commit and then on the lab
server; all 352 bug verdicts are identical. The HWASan formal run agrees with an earlier development
run in 95 of 96 cells, the exception being the tag collision above.

We then examined every one of the 174 misses from the raw output. In each, the program printed
\code{ST READY}, so the faulty line ran; 170 ended with \code{ST DONE} and exit 0, and the other four
are the case-4 crashes. No miss carries a report of any kind in its standard error --- we searched for
every ASan and HWASan report type, Memcheck's invalid accesses and frees, uninitialised uses, and the
UBSan and LowFat messages --- and every Memcheck miss ends with ``\code{ERROR SUMMARY: 0 errors}''.
In 140 of the 174, the same program is detected by the same tool on \code{malloc}, which shows that
the access is instrumented and executed; the rest are explained by case semantics (cases 10, 15 are
unsupported on \code{malloc}; case 3 under ASan and case 14 under HWASan are missed everywhere) or by
the tag collision. All misses fall into seven mechanisms:

\begin{center}\small
\begin{tabular}{lrl}
\toprule
mechanism & cells & cases \\
\midrule
the pool hides the object (unannotated layers) & 122 & 0--2, 5--11, 13--15 \\
the overflow lands in another live object & 20 & 3 \\
stale pointer to reissued storage (annotated, address-based) & 4 & 10 \\
invalid free without invalid access (poisoning, tagging) & 14 & 13, 14 \\
HWASan's 16-byte tag granularity (annotated) & 6 & 0--2 \\
past the pool, outside ASan's large mapping & 7 & 4 \\
HWASan tag collision & 1 & 4 \\
\bottomrule
\end{tabular}
\end{center}

None is caused by the build, by optimisation or by the harness. Two original predictions were
corrected by the measurements and are documented with the cases: case 3 under Memcheck on
\code{malloc} (detected, not missed) and case 4 under ASan on the pools (missed, not detected).

\subsection{HWASan on x86-64}
\label{sec:hwasan-x86}

HWASan on x86-64 normally needs LAM, which neither our lab VM nor the lab server's Xeon Gold 6230
provides; the default mode aborts at start-up. We also tried clang's experimental aliasing mode
(\code{-fsanitize-hwaddress-}\allowbreak\code{experimental-aliasing}), which maps the heap at several virtual addresses
so that a tagged pointer is still valid. It runs and detects both positive controls. Over the full
HWASan matrix it detected 25 bugs against 51 for the AArch64 run, and nothing at all inside Arena or
Free list, annotated or not. A probe explains why: the aliasing heap tags only allocations up to
64~KiB, with three-bit pointer tags; blocks of 256~KiB and more are mapped outside the alias region
with tag 0 and their accesses are not checked. Our 1~MiB pools are invisible to it, and an
annotation tag such as 17 no longer fits in the pointer. We therefore did not use this mode: the
tested x86-64 aliasing implementation cannot check accesses within our 1~MiB pools, even with
annotations. This result does not establish the same limitation for smaller pools or other modes.
